\documentclass[11pt,a4paper,sans]{moderncv}
\moderncvstyle{banking}
\moderncvcolor{blue}

% Force first/last name and section headings to render in moderncv blue (color1).
\renewcommand*{\firstnamestyle}[1]{{\fontsize{34}{36}\bfseries\upshape\color{color1}#1}}
\renewcommand*{\lastnamestyle}[1]{{\fontsize{34}{36}\bfseries\upshape\color{color1}#1}}
\renewcommand*{\sectionstyle}[1]{{\sectionfont\color{color1}#1}}

\usepackage{hyperref}
\hypersetup{
    colorlinks=true,
    linkcolor=blue,
    filecolor=magenta,
    urlcolor=blue,
    pdftitle={Walter Hernandez - CV},
    pdfpagemode=FullScreen,
}
\usepackage[scale=0.77]{geometry}
\usepackage{needspace}

\name{Walter}{Hernandez}
\address{Argentina}{}{}
\phone[mobile]{(+54) 11 2382-6514}
\email{walskp@gmail.com}
\extrainfo{\href{https://www.linkedin.com/in/wal-hernandez-dev}{LinkedIn}, \href{https://github.com/Wal-Hernandez}{GitHub}}

\begin{document}
\makecvtitle

\vspace{-4pt}

%------------------------------------------------------------------------------
\section{Perfil Profesional}
%------------------------------------------------------------------------------
\vspace{1pt}
Software Engineer SSR con experiencia en desarrollo full stack y fuerte orientaci\'{o}n hacia
backend y cloud. He participado en la migraci\'{o}n de un sistema legacy .NET 2.2 a Node.js + AWS,
desarrollando APIs REST y servicios backend con Lambda, EventBridge y SQS, y aprovisionando
infraestructura con AWS CDK. He integrado OpenAI APIs en flujos de procesamiento documental y
participado en la evoluci\'{o}n de arquitecturas hacia una mejor separaci\'{o}n de responsabilidades.
Aplico principios SOLID y patrones de dise\~{n}o en producci\'{o}n, reviso c\'{o}digo y trabajo en equipos
multidisciplinarios. Busco seguir desarroll\'{a}ndome en arquitectura cloud, buenas pr\'{a}cticas de
ingenier\'{i}a y equipos colaborativos con foco en calidad de software.

%------------------------------------------------------------------------------
\section{Experiencia Profesional}
%------------------------------------------------------------------------------
\vspace{1pt}

\needspace{5\baselineskip}
\cventry{Oct 2024--Presente}{Desarrollador Full Stack SSR}{La Naci\'{o}n}{Argentina}{H\'{i}brido/Remoto}{\vspace{3pt}
\begin{itemize}
\item Contribu\'{i} a la migraci\'{o}n de un sistema legacy .NET 2.2 de newsletters hacia Node.js + AWS,
desarrollando servicios backend orientados a eventos integrados mediante EventBridge y SQS,
aprovisionando infraestructura con AWS CDK. Configur\'{e} Dead Letter Queues para tolerancia a fallos
y monitoreo con CloudWatch.
\item Aplico principios SOLID y patrones de dise\~{n}o (Strategy, Factory, Repository, Dependency Injection)
en c\'{o}digo backend productivo y arquitectura de APIs.
\item Evolucion\'{e} de responsable del flujo de newsletters a desarrollador core en el monorepo del Paywall,
contribuyendo a funcionalidades de negocio y trabajando con equipos de Arquitectura, Base de Datos y
Frontend en el refinamiento t\'{e}cnico de tareas complejas.
\item Desplegu\'{e} servicios mediante pipelines CI/CD en Azure DevOps y realic\'{e} code reviews. Depur\'{e}
problemas en producci\'{o}n trazando eventos a trav\'{e}s de EventBridge, logs de ejecuci\'{o}n de Lambda
y CloudWatch.
\end{itemize}}

\vspace{3pt}
\cventry{Jul 2023--Oct 2024}{Desarrollador de Software Full Stack}{Eppical}{Argentina}{Remoto}{\vspace{3pt}
\begin{itemize}
\item Desarroll\'{e} integraciones backend para flujos de procesamiento documental con IA,
enrutando documentos a trav\'{e}s de OpenAI APIs y procesando resultados en servicios C\#/.NET y
Node.js consumidos por frontends Angular.
\item Particip\'{e} en la evoluci\'{o}n de la arquitectura del sistema, separando responsabilidades
entre frontend, backend y base de datos para mejorar la escalabilidad, el reuso de c\'{o}digo y
la mantenibilidad.
\item Desarroll\'{e} funcionalidades full-stack integrando frontends Angular con backends C\#/.NET y Node.js,
incluyendo componentes de UI, endpoints de API e interacciones con SQL Server. Refactoric\'{e} c\'{o}digo
aplicando buenas pr\'{a}cticas y particip\'{e} en discusiones t\'{e}cnicas para evaluar implementaciones.
\end{itemize}}

\vspace{3pt}
\cventry{Feb 2021--Jul 2023}{Desarrollador Full Stack}{SYNAgro}{Argentina}{Presencial/Remoto}{\vspace{3pt}
\begin{itemize}
\item Desarroll\'{e} aplicaciones de negocio para el sector agroindustrial desarrollando servicios backend en
C\#/.NET, interfaces en React e Ionic, y cambios en base de datos SQL.
\item Constru\'{i} interfaces web y mobile responsivas con React, TypeScript e Ionic, resolv\'{i} bugs en
producci\'{o}n e implement\'{e} funcionalidades orientadas al cliente.
\end{itemize}}

%------------------------------------------------------------------------------
\section{Competencias Clave}
%------------------------------------------------------------------------------
\vspace{1pt}
\begin{itemize}
\item \textbf{Backend y Arquitectura:} C\#, .NET, Node.js, TypeScript, dise\~{n}o de APIs REST,
arquitectura event-driven, microservicios backend, principios SOLID, patrones de dise\~{n}o
(Strategy, Factory, Repository, DI), modernizaci\'{o}n de sistemas legacy.
\item \textbf{Cloud y DevOps:} AWS (Lambda, API Gateway, EventBridge, SQS, CloudWatch), AWS CDK,
Azure DevOps, pipelines CI/CD, Git Flow.
\item \textbf{Integraci\'{o}n con IA:} OpenAI APIs, pipelines de procesamiento documental,
desarrollo de funcionalidades con IA.
\item \textbf{Bases de Datos:} SQL Server, PostgreSQL, MongoDB.
\item \textbf{Frontend:} React, Angular, Ionic.
\item \textbf{Pr\'{a}cticas de Ingenier\'{i}a:} Code reviews, colaboraci\'{o}n cross-functional,
documentaci\'{o}n t\'{e}cnica, trabajo con Product Owners y equipos de negocio.
\end{itemize}

%------------------------------------------------------------------------------
\section{Educaci\'{o}n}
%------------------------------------------------------------------------------
\vspace{1pt}

\needspace{5\baselineskip}
\cventry{2024--Presente}{Seguridad Inform\'{a}tica (En Curso)}{Universidad Siglo 21}{Argentina (Remoto)}{}{\vspace{1pt}
Enfoque: infraestructura, seguridad y fundamentos de tecnolog\'{i}as modernas.
}

\vspace{3pt}
\cventry{2021--2024}{Ingenier\'{i}a en Sistemas de Informaci\'{o}n (Parcial, Reorientada)}{Universidad Tecnol\'{o}gica Nacional (UTN)}{Argentina}{}{\vspace{1pt}
Materias relevantes: Matem\'{a}tica, L\'{o}gica, Estructuras de Datos, Algoritmos, Ingenier\'{i}a de Software.
}

\vspace{3pt}
\cventry{2021--2022}{Bootcamp de Desarrollo Web Full Stack}{Henry}{Remoto}{}{\vspace{1pt}
M\'{a}s de 1500 horas de formaci\'{o}n intensiva, pair programming diario, proyectos individuales y grupales.
}

%------------------------------------------------------------------------------
\section{Idiomas}
%------------------------------------------------------------------------------
\vspace{1pt}
\begin{itemize}
\item \textbf{Espa\~{n}ol:} Nativo
\item \textbf{Ingl\'{e}s:} C1 Avanzado
\end{itemize}

\end{document}
